\par\noindent\textbf{(1)} 耗尽层p侧电荷 \(\rho=-qN_A\)、n侧 \(\rho=qN_D\)，Poisson方程 \(dE/dx=\rho/\epsilon_s\)，边界 \(E(-x_p)=E(x_n)=0\)。界面电场连续要求 \(N_Ax_p=N_Dx_n\)，故 \(x_p=N_DW/(N_A+N_D)\)、\(x_n=N_AW/(N_A+N_D)\)，\(W=x_p+x_n\)，是几何电荷中性而非电子与空穴浓度相等。
\par\noindent\textbf{(2)} 静电势差是两段电场三角形积分：\(V_{bi}-V=q(N_Ax_p^2+N_Dx_n^2)/(2\epsilon_s)\)。用 \(N_Ax_p=N_Dx_n\) 消去两未知量，得到 \(W^2=(2\epsilon_s/q)(1/N_A+1/N_D)(V_{bi}-V)\)。因此强反偏使耗尽宽度增大；正偏逼近内建势时耗尽近似最终失效。
\par\noindent\textbf{(3)} 单位面积耗尽电荷模 \(Q_A=qN_Ax_p=qN_AN_DW/(N_A+N_D)\)。微分单位面积结电容 \(C_j/A=|dQ_A/dV|=\epsilon_s/W\)，所以 \(1/C_j^2=[2(V_{bi}-V)/(q\epsilon_sA^2)](1/N_A+1/N_D)\)。测量反偏下 \(1/C_j^2\) 对电压的斜率可反演掺杂，但需注意串联电阻和界面陷阱。
\par\noindent\textbf{(4)} 在低注入条件下少子扩散满足 \(D_n\delta n''-\delta n/\tau_n=0\)，解 \(\delta n(x)=\delta n(0)e^{-x/L_n}\)，\(L_n=\sqrt{D_n\tau_n}\)。结边界 \(\delta n(0)=n_{p0}(e^{qV/(k_BT)}-1)\)，电子、空穴扩散电流相加为 \(I=I_s(e^{qV/(k_BT)}-1)\)，\(I_s=qA(D_nn_{p0}/L_n+D_pp_{n0}/L_p)\)，实际大注入、复合和电阻使其偏离。
